\section{Open Questions and Next Steps}

The replication surfaced five substantive open questions. Each is grounded in the paper's own gaps and has a concrete, free-tool next step.

\subsection{Q1: BER-polymerase specificity}
Does the radiosensitization phenotype extend to POL$\lambda$ / POL$\mu$ when knocked down in the same PA1 background, or is it specific to POL$\beta$?

\textbf{Basis.} The paper attributes PA1Pol$\beta\Delta$ radiosensitivity to loss of gap-filling BER capacity, but POL$\lambda$ and POL$\mu$ also contribute to BER (and to NHEJ). If the phenotype is BER-mediated, isogenic KDs of POL$\lambda$ / POL$\mu$ should produce a similar DMF shift. If PA1Pol$\beta\Delta$ is uniquely radiosensitive, the mechanism may be dominant-negative rather than pure BER-loss.

\textbf{Next steps.} Design isogenic PA1 KD panel (siRNA or CRISPRi against POLB, POLL, POLM, plus double-KDs). Score clonogenic survival at 0/5/10/15~Gy with $n \geq 6$. Refit LQ per line and compare DMFs. In silico prerequisite: repeat our sequence/PDB audit for POL$\lambda$ (PDB 2BCQ) and POL$\mu$ (PDB 2IHM).

\subsection{Q2: Interaction with DSB-repair models}
How does the reported PA1Pol$\beta\Delta$ radiosensitivity interact with modern DSB-repair models (HR / NHEJ / MMEJ) that the paper does not touch?

\textbf{Basis.} The mutant LQ fit gives $\alpha/\beta = 0.30$~Gy (essentially no shoulder), more consistent with catastrophic DSB-repair failure than BER-only loss. BER handles SSBs / abasic sites; SSBs converted to DSBs at replication forks would recruit HR/NHEJ. A pure BER defect should not collapse the shoulder that dramatically.

\textbf{Next steps.} Integrate the mutant LQ parameters into a mechanistic model (MEDRAS, LEM-IV, or a lightweight two-lesion Markov model) with explicit BER, HR, NHEJ compartments. Fit BER-only vs.\ BER+HR-crosstalk hypotheses to the shoulder shape. Free tool: MEDRAS Python port on GitHub. Predict SF at 2/6/8~Gy for validation.

\subsection{Q3: Tumor-genotype context}
Which tumor genotypes with pre-existing BER deficiency (MUTYH-associated polyposis, OGG1-deficient, XRCC1 hypomorphs) are candidate clinical contexts where POL$\beta$ targeting would compound radiosensitization?

\textbf{Basis.} If loss of a single BER component (POL$\beta$) produces DMF $\sim$1.8, combining POL$\beta$ inhibition with upstream BER defects should be synthetically lethal under radiation. The paper hints at this but does not explore. Public cancer genomics (TCGA, DepMap) can map POLB expression against BER-pathway alteration burden.

\textbf{Next steps.} Query DepMap for lines with POLB LOF or low expression cross-tabulated with MUTYH / OGG1 / XRCC1 status. Use CellMinerCDB to correlate POLB expression with radiation sensitivity (AUC). Extract cancer types where the combined-deficiency profile is enriched. Free endpoints: DepMap portal, cBioPortal, CellMinerCDB.

\subsection{Q4: Template / isoform sensitivity of the docking story}
How sensitive is the docking / structural story to POL$\beta$ isoform choice (canonical 335-aa vs.\ splice variants) and to the specific PDB template used (1TV9 vs.\ 1BPY, 3ISB, 4KLE)?

\textbf{Basis.} The paper uses only 1TV9 (Pol$\beta$ bound to nicked mismatched DNA). Human POLB has splice isoforms and multiple crystal structures with different DNA substrates and conformations (open vs.\ closed thumb). Docking scores against BER partners depend on receptor conformation; a single-template analysis may bias the panel.

\textbf{Next steps.} Repeat the PDB structural audit for $\geq$3 additional Pol$\beta$ templates (1BPY apo, 3ISB substrate complex, 4KLE with Mn). Re-run top-3 ClusPro dockings (NEIL1, XRCC1, APE1) against each variant. Compare rank stability. Free endpoints: RCSB, ClusPro (single-user quota permitting).

\subsection{Q5: ML-driven radiosensitizer prediction}
Can machine-learning models trained on BER-pathway expression signatures predict radiosensitizer response in patient tumors, using PA1Pol$\beta\Delta$ as a positive training example?

\textbf{Basis.} The paper generates a single-line phenotype. Modern pan-cancer datasets (CCLE, GDSC, DepMap PRISM) include radiation response for hundreds of lines with paired transcriptomics. A simple GBM or logistic model on BER-pathway genes may recover a signature that generalizes across tumor types.

\textbf{Next steps.} Assemble training data from DepMap (radiation AUC) + CCLE (RNA-seq). Feature set: 10 BER genes (POLB, OGG1, NEIL1, XRCC1, PNKP, APE1, PARP1, PARP2, FEN1, LIG3) + POLB LOF status. Model: L1-regularized logistic regression, 5-fold CV. Validate on TCGA cohorts with radiation-therapy outcome. Free endpoints: DepMap portal, TCGA-GDC, PORTAL Kaplan-Meier plotter.
